\section{Deterministic Finite Automata}

Literature:
\begin{itemize}
    \item \cite{crosswhite_finite_2008} Basic Representation by TT/MPS.
    \item \cite{bellante_compiling_2026} Usage of canonical formats (orthogonalization) to find quantum circuit representations.
\end{itemize}

Idea:
\begin{itemize}
    \item Construct chain graphs decorated with basis encodings of state transition functions, as in \cite{crosswhite_finite_2008,bellante_compiling_2026}
    \item Contract with basis encoding of halting states and get the boolean tensor indicating the accepted states of the automaton
\end{itemize}

\begin{definition}[Deterministic finite automaton (DFA), see Def~2 in \cite{bellante_compiling_2026}]
    A tuple of
    \begin{itemize}
        \item set of internal states $Q$
        \item alphabet $\Sigma$
        \item initial state $r_0\in Q$
        \item accepted states $F\subset Q$
        \item transition function $\delta:Q\otimes \Sigma\rightarrow Q\cup\{\varnothing\}$
    \end{itemize}
    is called a deterministic finite automaton (DFA).
\end{definition}

We enumerate the set $Q$ by an map $\indexinterpretationof{Q}$ with enumeration variable $\indvariableof{Q}$, and the alphabet $\Sigma$ by a map $\indexinterpretationof{\Sigma}$ with an enumeration variable $\indvariableof{\Sigma}$.
Both indicator variables will be copied for $\catenumeratorin$ (respectively $\catenumerator\in[\catorder+1]$) and the copies denoted by $\indvariableof{\Sigma,\catenumerator}$ and $\indvariableof{Q,\catenumerator}$.
The acceptance tensor of order $\catorder\in\nn$ is the boolean tensor $\hypercoreat{\indvariableof{\Sigma,[\catorder]}}$ indicating whether the automaton accepts the input sequence of length at most $\catorder$.
Its coordinates are
\begin{align*}
    \hypercoreat{\indexedindvariableof{\Sigma,[\catorder]}}
    = \begin{cases}
          1 & \ifspace \text{The automaton accepts the input sequence }\left(\indexinterpretationofat{\Sigma}{\indexedindvariableof{\catenumerator}}\wcols\catenumeratorin\right)\\
          0 & \elsetext
    \end{cases} \, .
\end{align*}


\begin{lemma}[Tensor representation]
    Let $(Q,\Sigma,\delta,r_0,F)$ be a deterministic finite automaton.
    Then the acceptance tensor of order $\catorder\in\nn$ is the contraction
    \begin{align*}
        \hypercoreat{\indvariableof{\Sigma,[\catorder]}}
        = \contractionof{
            \left\{ \bencodingofat{\delta}{\indvariableof{Q,\catenumerator},\indvariableof{\Sigma,\catenumerator},\indvariableof{\Sigma,\catenumerator}} \wcols \catenumeratorin
            \right\}
            \cup
            \{\onehotmapofat{\invindexinterpretationofat{Q}{r_0}}{\indvariableof{Q,0}},
            \onehotmapofat{\invindexinterpretationofat{Q}{F}}{\indvariableof{Q,\catorder}}
            \}
        }{\indvariableof{\Sigma,[\catorder]}} \, .
    \end{align*}
\end{lemma}
\begin{proof}
    By basis calculus:
    Given the one-hot encoding of the input word and contracting with the tensor network (expect the basis encoding tensor of $F$) leaving $\indvariableof{Q,\catorder}$ open, we get the one hot encoding of the finial state of the automaton.
    Contracting with the one-hot encoding of $F$ is then $1$, if and only if the final state is in the accepted states.
\end{proof}

% Complement
Note that the acceptance tensor of the complement DFA is prepared when replacing the basis encoding $\onehotmapofat{\invindexinterpretationofat{Q}{F}}{\indvariableof{Q,\catorder}}$ by the basis encoding $\onehotmapofat{\invindexinterpretationofat{Q}{\bar{F}}}{\indvariableof{Q,\catorder}}$ of the complement of the accepting states.


\subsection{Outlook}

We can use this theory to:
\begin{itemize}
    \item Criterion that a language is not regular: Unbounded matrification ranks in the word tensors.
    For example, this shows that $\{0^n1^n \wcols n\in\nn\}$ is not regular.
    \item Proof that regular languages are closed under regular operations: Use parallel TT cores, modifications in the transition functions.
\end{itemize}